
\bigskip
이상의 교차분석과 논의를 종합하면, $\DRTO$ 개별 시그모이드 바운딩 모델은
관측성과 불확실성의 효과를 구조적으로 분리하고, 극단 영역에서의
비선형 상호작용을 이중채널 감쇠로 포착하며,
전문가 검증에 의해 실무적 타당성이 확인된 프레임워크이다.
이어지는 절에서는 이러한 이론적 발견이 BCM 실무와 정책 수립에 갖는
구체적 시사점을 논의한다.


%% ===================================================================
%% 이하 구 제6장(시사점) 내용 — ch5와 병합
%% ===================================================================

이하에서는 \chref{ch:results}의 시뮬레이션 실험과 앞서 논의한 교차 조건 분석 결과를
토대로 $\DRTO$ 프레임워크의 시사점을 이론적, 방법론적, 그리고 실무적 관점에서
체계적으로 논의한다. 먼저 개별 시그모이드 바운딩 모델이 BCM 이론에
기여하는 바를 정리하고(\secref{sec:theoretical-implications}),
연구 방법론의 혁신과 확장 가능성을 논의한다(\secref{sec:methodological-implications}).
이어서 2단계 프레임워크, 에스컬레이션 체계, 6개 산업 적용 가이드라인 등
실무적 시사점을 제시하고(\secref{sec:practical-implications}),
마지막으로 연구의 한계점을 투명하게 논의한다(\secref{sec:limitations}).
